\section{Visualising forecasts}
\label{sec:visual}

The right chart depends on who reads it and which decision follows (Table~\ref{tab:viz}). Four rules apply everywhere. Show uncertainty, as quantile bands or ensemble members, rather than a single line. State the time zone and interval convention and mark the present and the relevant gate closures. Use one categorical colour per entity throughout a report, never two y-axes in one chart, and dashed lines only for references and thresholds. Publish the numbers behind every chart, so that the chart is never the only way to read a value.

The figures in this paper follow these rules with a colour palette validated for colour-vision deficiency: neighbouring categorical colours are separated by at least 9 units of $\Delta E$ (OKLab$\times$100) under simulated protanopia and deuteranopia and by at least 19.6 for normal vision, magnitudes use a single-hue blue ramp, and every figure's data are in \texttt{data/processed}.

% Table 8: chart types by audience.
\begin{table}[t]
\centering\small
\caption{\textbf{Charts that answer each audience's question.}}
\label{tab:viz}
\setlength{\tabcolsep}{3pt}
\begin{tabular}{@{}L{1.85cm}L{5.95cm}@{}}
\toprule
Audience & Charts \\
\midrule
Trading & fan chart of quantiles with actuals; run-to-run change of the forecast; ensemble plume; residual load with price forecast; heat map of quarter-hours by day \\
Operations & meteogram with activity limits (gusts, waves, lightning); warning maps with plant locations; satellite loops \\
Asset management & expected vs.\ actual; power vs.\ irradiance or wind scatter; carpet plot of day by time of day; loss waterfall; cumulative yield vs.\ budget with P50/P90 band \\
Development & resource maps; wind rose; Weibull histogram; month-by-hour profile; exceedance curve; interannual variability \\
Construction & weather-window calendar; monthly workability from hindcasts \\
Forecast analysts & error and skill by lead time; reliability diagram; rank and PIT histograms; error by hour and month; worst-days gallery \\
\bottomrule
\end{tabular}
\end{table}

